\documentclass[12pt,a4paper]{article}
\usepackage[utf8]{inputenc}
\usepackage{amsmath,amsfonts,amssymb}
\usepackage{listings}
\usepackage{xcolor}
\usepackage{geometry}
\usepackage{hyperref}
\usepackage{fancyhdr}

\geometry{margin=1in}
\pagestyle{fancy}
\fancyhf{}
\rhead{NumPy - Complete Guide}
\lfoot{Python for Data Science}
\rfoot{\thepage}

\lstset{
    language=Python,
    basicstyle=\ttfamily\small,
    keywordstyle=\color{blue}\bfseries,
    commentstyle=\color{green!60!black},
    stringstyle=\color{red},
    numbers=left,
    numberstyle=\tiny\color{gray},
    stepnumber=1,
    numbersep=8pt,
    showstringspaces=false,
    breaklines=true,
    frame=single,
    backgroundcolor=\color{gray!10}
}

\title{\textbf{NumPy - Complete Guide}}
\author{Python for Data Science - Beginner's Level}
\date{\today}

\begin{document}

\maketitle
\tableofcontents
\newpage

\section{Overview}

NumPy (Numerical Python) is the fundamental package for scientific computing in Python. It provides powerful N-dimensional array objects and tools for working with arrays.

\section{Installation and Import}

\begin{lstlisting}
# Installation
# pip install numpy

# Import
import numpy as np
\end{lstlisting}

\section{Array Creation}

\subsection{Basic Array Creation}

\begin{lstlisting}
# From lists
arr1d = np.array([1, 2, 3, 4, 5])
arr2d = np.array([[1, 2, 3], [4, 5, 6]])
arr3d = np.array([[[1, 2], [3, 4]], [[5, 6], [7, 8]]])

# Array properties
print(arr2d.shape)    # (2, 3)
print(arr2d.dtype)    # int64
print(arr2d.ndim)     # 2
print(arr2d.size)     # 6
\end{lstlisting}

\subsection{Specialized Array Creation}

\begin{lstlisting}
# Zeros and ones
zeros = np.zeros((3, 4))           # 3x4 array of zeros
ones = np.ones((2, 3))             # 2x3 array of ones
full = np.full((2, 2), 7)          # 2x2 array filled with 7

# Identity and diagonal
identity = np.eye(3)               # 3x3 identity matrix
diag = np.diag([1, 2, 3])         # Diagonal matrix

# Range arrays
arange = np.arange(0, 10, 2)       # [0, 2, 4, 6, 8]
linspace = np.linspace(0, 1, 5)    # [0, 0.25, 0.5, 0.75, 1]

# Random arrays
random = np.random.rand(3, 3)      # Random values [0, 1)
randint = np.random.randint(1, 10, (2, 3))  # Random integers
\end{lstlisting}

\section{Array Indexing and Slicing}

\subsection{Basic Indexing}

\begin{lstlisting}
arr = np.array([[1, 2, 3, 4],
                [5, 6, 7, 8],
                [9, 10, 11, 12]])

# Single element
print(arr[1, 2])      # 7

# Row and column
print(arr[1, :])      # [5 6 7 8] (row 1)
print(arr[:, 2])      # [3 7 11] (column 2)

# Slicing
print(arr[0:2, 1:3])  # [[2 3], [6 7]]
\end{lstlisting}

\subsection{Advanced Indexing}

\begin{lstlisting}
arr = np.array([1, 2, 3, 4, 5, 6, 7, 8, 9, 10])

# Boolean indexing
mask = arr > 5
print(arr[mask])      # [6 7 8 9 10]

# Fancy indexing
indices = [1, 3, 5]
print(arr[indices])   # [2 4 6]

# Conditional selection
result = np.where(arr > 5, arr, 0)  # Replace values ≤5 with 0
\end{lstlisting}

\section{Array Operations}

\subsection{Mathematical Operations}

\begin{lstlisting}
arr1 = np.array([1, 2, 3, 4])
arr2 = np.array([5, 6, 7, 8])

# Element-wise operations
print(arr1 + arr2)    # [6 8 10 12]
print(arr1 * arr2)    # [5 12 21 32]
print(arr1 ** 2)      # [1 4 9 16]

# Universal functions
print(np.sqrt(arr1))  # [1. 1.414 1.732 2.]
print(np.exp(arr1))   # [2.718 7.389 20.086 54.598]
print(np.sin(arr1))   # [0.841 0.909 0.141 -0.757]
\end{lstlisting}

\subsection{Aggregation Functions}

\begin{lstlisting}
arr = np.array([[1, 2, 3],
                [4, 5, 6]])

# Basic aggregations
print(np.sum(arr))        # 21
print(np.mean(arr))       # 3.5
print(np.std(arr))        # 1.707
print(np.min(arr))        # 1
print(np.max(arr))        # 6

# Axis-specific aggregations
print(np.sum(arr, axis=0))    # [5 7 9] (column sums)
print(np.sum(arr, axis=1))    # [6 15] (row sums)
\end{lstlisting}

\section{Array Manipulation}

\subsection{Reshaping}

\begin{lstlisting}
arr = np.arange(12)  # [0 1 2 3 4 5 6 7 8 9 10 11]

# Reshape
reshaped = arr.reshape(3, 4)

# Flatten
flattened = reshaped.flatten()

# Transpose
transposed = reshaped.T
\end{lstlisting}

\subsection{Joining and Splitting}

\begin{lstlisting}
arr1 = np.array([[1, 2], [3, 4]])
arr2 = np.array([[5, 6], [7, 8]])

# Concatenation
hstack = np.hstack([arr1, arr2])  # [[1 2 5 6], [3 4 7 8]]
vstack = np.vstack([arr1, arr2])  # [[1 2], [3 4], [5 6], [7 8]]

# Splitting
arr = np.arange(8)
split = np.split(arr, 4)  # [array([0, 1]), array([2, 3]), ...]
\end{lstlisting}

\section{Linear Algebra}

\begin{lstlisting}
# Matrix operations
A = np.array([[1, 2], [3, 4]])
B = np.array([[5, 6], [7, 8]])

# Matrix multiplication
dot_product = np.dot(A, B)    # or A @ B

# Matrix properties
det = np.linalg.det(A)        # Determinant
inv = np.linalg.inv(A)        # Inverse
eigenvals = np.linalg.eigvals(A)  # Eigenvalues

# Solving linear systems Ax = b
b = np.array([1, 2])
x = np.linalg.solve(A, b)
\end{lstlisting}

\section{Broadcasting}

\begin{lstlisting}
# Broadcasting allows operations between arrays of different shapes
arr = np.array([[1, 2, 3],
                [4, 5, 6]])

# Scalar broadcasting
result1 = arr + 10            # Add 10 to all elements

# Vector broadcasting
vec = np.array([10, 20, 30])
result2 = arr + vec           # Add vector to each row

# Column vector broadcasting
col_vec = np.array([[10], [20]])
result3 = arr + col_vec       # Add column vector
\end{lstlisting}

\section{Performance Tips}

\begin{lstlisting}
# Use vectorized operations instead of loops
# Slow
result_slow = []
for i in range(len(arr)):
    result_slow.append(arr[i] ** 2)

# Fast
result_fast = arr ** 2

# Use appropriate data types
int_arr = np.array([1, 2, 3], dtype=np.int32)    # Less memory
float_arr = np.array([1, 2, 3], dtype=np.float64) # More precision
\end{lstlisting}

\section{Summary}

NumPy provides:
\begin{itemize}
    \item \textbf{Efficient arrays}: Fast numerical operations
    \item \textbf{Broadcasting}: Operations between different shapes
    \item \textbf{Linear algebra}: Matrix operations and decompositions
    \item \textbf{Random numbers}: Statistical distributions and sampling
    \item \textbf{Universal functions}: Element-wise operations
    \item \textbf{Memory efficiency}: Optimized C implementations
\end{itemize}

Essential for scientific computing and data science in Python!

\end{document}